\section{源代码与复现说明}
\subsection{源码组织}
源码位于项目的 \code{pq-test/} 下。以下给出与本报告对应的文件，
完整源码是否需要逐页附入报告，以当前课程要求为准。
\begin{table}[H]\centering\small
\caption{主要源码与功能对应}
\begin{tabularx}{\linewidth}{lX}\toprule
文件 & 功能\\\midrule
\code{csr-graph.hpp}、\code{dimacs-parser.hpp} & 共享图结构与文件解析\\
\code{general-dijkstra.hpp} & 共用 Dijkstra 主循环\\
\code{heap-adaptor.hpp} & 时间戳接口、标准库与 PBDS 适配器\\
\code{binary-heap.hpp}、\code{segment-tree-heap.hpp} & 二叉堆与段树\\
\code{fibonacci-heap.hpp}、\code{radix-heap.hpp} & Fibonacci 与 radix heap\\
\code{general-bellman-ford.hpp} & 两种标签修正算法\\
\code{dij-bench-v2.cpp} & 随机图参数扫描\\
\code{dij-bench-complexity.cpp} & 固定密度复杂度实验\\
\code{dij-bench-dimacs.cpp} & 路网查询批次基准\\
\code{heap-adaptor-test.cpp}、\code{dimacs-parser-test.cpp} & 接口、差分、回归及解析测试\\\bottomrule
\end{tabularx}\end{table}

下列节选保留实际主循环逻辑，省略文件头中的身份信息，并补充解释性注释。
它用于说明共享接口，不替代完整源码提交。

\noindent\begin{minipage}{\linewidth}
\lstinputlisting[language=C++,caption={共享主循环节选}]{code/dijkstra.hpp}
\end{minipage}

\drafttodo{课程结构要求规定提交源码至少 30\% 的行带有注释。此底稿没有声称项目全部源码已满足该比例。
正式提交前需按课程口径审查实际提交的源码，并补全必要注释；不能用本页注释比例替代整体审查。}

\subsection{数据与作图复现}
\code{reports/data/} 保留 CSV 快照、测量校验结果及经过匿名化处理的溯源信息。

总表 \code{all-results.csv} 含 2151 行；计时汇总
\code{timing-summary.csv} 含 315 组配置。
脚本从 \code{run_type=iteration} 的 909 条原始测量行重新计算均值和标准差，
不会把均值、CV 或 RMS 行当作新样本。

在项目根目录执行：
\begin{lstlisting}
python -m pip install -r tools/requirements-plot.txt
python tools/plot_benchmarks.py
python reports/build_report.py
\end{lstlisting}
各作图脚本也可独立运行，并通过 \code{--data-dir}、\code{--output-dir} 指定目录。
图同时保存为矢量 PDF 和 PNG；表格由 \code{generate_report_tables.py} 生成到
\code{reports/tables/}，不应手工改动数值。

从源码重新执行 benchmark 的基本步骤如下；不同机器请自行指定构建目录和路网数据目录。

\noindent\begin{minipage}{\linewidth}
\begin{lstlisting}
cmake -S . -B build-wsl -G Ninja -DCMAKE_BUILD_TYPE=Release
cmake --build build-wsl --target heap-adaptor-test dimacs-parser-test \
  dij-bench-v2 dij-bench-complexity dij-bench-dimacs
ctest --test-dir build-wsl --output-on-failure
build-wsl/pq-test/dij-bench-v2 --benchmark_min_time=0.5s \
  --benchmark_repetitions=3 --benchmark_out=random.json \
  --benchmark_out_format=json
build-wsl/pq-test/dij-bench-complexity --benchmark_min_time=0.5s \
  --benchmark_repetitions=3 --benchmark_out=complexity.json \
  --benchmark_out_format=json
build-wsl/pq-test/dij-bench-dimacs --gr DATA/USA-road-d.NY.gr \
  --benchmark_repetitions=1 --benchmark_out=ny.json \
  --benchmark_out_format=json
build-wsl/pq-test/dij-bench-dimacs --gr DATA/USA-road-d.BAY.gr \
  --benchmark_repetitions=1 --benchmark_out=bay.json \
  --benchmark_out_format=json
\end{lstlisting}
\end{minipage}
本地原始运行目录另保留从 JSON 导出 CSV 的脚本和完整命令。
若采用新的参数或查询集，应保留新的溯源文件并重新生成图表；本报告数据快照不会自动替换。
\FloatBarrier
